% ── Rozwiązanie: Zachłanny wybór zajęć -- instancja 1 ─────────────
\subsection{Zachłanny wybór zajęć -- instancja 1}

Dane ($n = 7$ zajęć), już posortowane po~rosnącym czasie zakończenia $t_i$.
Algorytm \textsc{GreedyActivitySelector} wybiera $z_1$, a~następnie idzie po
zajęciach rosnąco i~dodaje $z_i$ do~rozwiązania $A$ wtedy i~tylko wtedy, gdy jest
zgodne z~ostatnio wybranym $z_j$, tj. $s_i \geq t_j$. Wybrane zajęcia wyróżniono
na~\textcolor{accentB}{pomarańczowo}.

\begin{aztable}
	\centering
	\renewcommand{\arraystretch}{1.2}
	\begin{NiceTabular}{c|ccccccc}[margin,hvlines]
		$i$   & \color{accentB}1 & 2 & 3 & \color{accentB}4 & 5 & 6 & \color{accentB}7  \\
		\hline
		$s_i$ & \color{accentB}1 & 3 & 0 & \color{accentB}5 & 3 & 5 & \color{accentB}8  \\
		$t_i$ & \color{accentB}4 & 5 & 6 & \color{accentB}7 & 9 & 9 & \color{accentB}10
	\end{NiceTabular}
\end{aztable}

\begin{dygresja}
	Przebieg pętli ($j$ -- indeks ostatnio wybranego zajęcia, początkowo
	$A = \{z_1\}$, $j = 1$, $t_j = 4$):
	\begin{itemize}
		\item $i = 2$: $s_2 = 3 \geq t_j = 4$? \emph{nie} -- pomijamy.
		\item $i = 3$: $s_3 = 0 \geq 4$? \emph{nie} -- pomijamy.
		\item $i = 4$: $s_4 = 5 \geq 4$? \emph{tak} -- $A \cup \{z_4\}$, $j \Assign 4$, $t_j = 7$.
		\item $i = 5$: $s_5 = 3 \geq 7$? \emph{nie} -- pomijamy.
		\item $i = 6$: $s_6 = 5 \geq 7$? \emph{nie} -- pomijamy.
		\item $i = 7$: $s_7 = 8 \geq 7$? \emph{tak} -- $A \cup \{z_7\}$, $j \Assign 7$.
	\end{itemize}
\end{dygresja}

\paragraph{Wynik.}
Wybrane zajęcia są parami zgodne, bo ich przedziały
$\langle 1, 4)$, $\langle 5, 7)$, $\langle 8, 10)$ są rozłączne:
\[
	\boxed{A = \{z_1,\ z_4,\ z_7\}}
\]
o~liczności $|A| = 3$. 
